\section{Weisfeiler--Leman conventions and the bijective game}
\label{game:section}

The joint refinement and its $(k+1)$-pair characterization are standard;
see Kiefer and Neuen \cite[Section~2.2 and Theorem~1]{KieferNeuen2019}.
We prove the correspondence here, including the separate convention,
to fix both the parameter shift and the order of the game moves.

Graphs in this section are finite, simple, and nonempty.  They may carry
vertex colors from a common palette; isomorphisms must preserve these
colors.  The uncolored case uses a single vertex color.  For an ordered
$k$-tuple $\bar v$ and a vertex $z$, write $\bar v[i\leftarrow z]$ for
replacement of coordinate $i$ by $z$.

\begin{definition}[The two refinement conventions]
\label{game:conventions}
Let $k\geq 2$.  The initial color $\chi_0^X(\bar v)$ records the atomic
type of $\bar v$ in $X$: all coordinate equalities, adjacencies, and vertex
colors.  The \emph{joint} convention updates by
\[
 \chi_{t+1}^X(\bar v)=
 \left(\chi_t^X(\bar v),
 \left\{\!\left\{
   \bigl(\chi_t^X(\bar v[i\leftarrow z])\bigr)_{i=1}^k
       :z\in V(X)
 \right\}\!\right\}\right).
\]
The \emph{separate} convention instead records the ordered list
\[
 \left(
 \left\{\!\left\{
   \chi_t^X(\bar v[i\leftarrow z]):z\in V(X)
 \right\}\!\right\}\right)_{i=1}^k
\]
along with the old color.  In either convention, the same recursive color
names are used for both input graphs.  Two graphs are equivalent if their
$k$-tuple color histograms agree at every round.
\end{definition}

The partitions of $V(G)^k\sqcup V(H)^k$ induced by these common color
names refine monotonically and eventually stop splitting.  Write
$\bar u\sim\bar v$ for equality of the resulting stable colors, with
the convention understood.  Thus stable histogram equality is equivalent
to equality at every round.  For the joint convention, stability gives
the following property: if $\bar u\sim\bar v$, there is a bijection
$f:V(G)\to V(H)$ such that
\begin{equation}
 \bar u[i\leftarrow z]\sim\bar v[i\leftarrow f(z)]
 \quad\text{for every }z\in V(G)\text{ and every }i\in\{1,\ldots,k\}.
 \label{game:joint-extension}
\end{equation}
Indeed, match equal members of the two finite multisets of replacement
color vectors.  For the separate convention, the same conclusion holds
for a specified coordinate $i$, with a bijection $f_i$ allowed to depend
on $i$.  An induction on the round also shows that $\sim$ is preserved
by simultaneous permutations of the coordinates.

\begin{definition}[Bijective game]
\label{game:definition}
The bijective game with $K$ pairs is played on $G,H$ as follows.  Spoiler
may remove any marked pairs and choose a vacant pair.  Duplicator then
announces a bijection $f:V(G)\to V(H)$, after which Spoiler chooses
$z\in V(G)$ and marks $z,f(z)$ with that pair.  Duplicator loses as soon
as the marked correspondence fails to preserve equality, adjacency, or
vertex colors.  A winning strategy for Duplicator must survive every
finite continuation, and hence every infinite play.  If the graph orders
differ, Duplicator cannot announce a bijection and loses.
\end{definition}

A winning position remains winning after pairs are removed: Duplicator
may retain the removed pairs privately and use the original strategy,
forgetting a private pair if its label is reused. The public position
is always a restriction of the simulated winning position.

\begin{lemma}[Exact stable equivalence]
\label{game:wl}
For the joint convention put $K=k+1$, and for the separate convention put
$K=k$.  For graphs of equal positive order the following hold.
\begin{enumerate}
 \item Two $k$-tuples have the same stable color if and only if
 Duplicator has a winning $K$-pair strategy from the position marking
 those tuples on $k$ corresponding pairs.
 \item The graphs have equal stable $k$-tuple color histograms if and
 only if Duplicator has a winning $K$-pair strategy from the empty
 position.
\end{enumerate}
\end{lemma}

\begin{proof}
First use the joint convention.  Maintain full $K$-tuples whose every
$k$-coordinate projection has equal stable colors; entries in vacant
slots are virtual.  A pair of related $k$-tuples can be extended to such
full tuples by choosing any $z$ and using
\eqref{game:joint-extension}: the projection omitting the new slot is
the original tuple, and every other projection is a replacement, up to
a simultaneous coordinate permutation.

To fill a vacant slot, apply \eqref{game:joint-extension} to the other
$k$ slots.  The announced bijection preserves the invariant for every
choice by Spoiler.  Removing pebbles simply makes their entries virtual.
Because $k\geq2$, every pair of marked entries belongs to some
$k$-coordinate projection, so the invariant implies partial
isomorphism.  This proves the forward implication in the tuple
statement.  Equal histograms supply a related pair of $k$-tuples, from
which the same construction starts with all entries virtual.

Conversely, from any winning position with $k$ marked pairs, prove by
induction on $t$ that the two tuples have equal colors at round $t$.
The case $t=0$ is partial isomorphism.  For the inductive step, ask to
use the spare pair.  For every selection $z$, retain the new pair and
discard any one of the original pairs.  The resulting $k$-pair position
is still winning, and its round-$t$ colors agree by induction.
The same announced bijection works for all discarded coordinates;
therefore the joint replacement multisets agree.

For the separate convention maintain a related full $k$-tuple, again
using virtual entries in vacant slots.  Fill slot $i$ with the bijection
$f_i$ supplied by its separate replacement multiset.  The resulting full
tuple remains related, and hence gives partial isomorphism.  For the
converse induction, remove pair $i$ and replace it using the winning
strategy.  This gives exactly the equality of replacement multisets
needed in coordinate $i$.

Finally, from an empty-position winning strategy, place the $k$ pairs
in a fixed order.  The successive announced bijections induce a
bijection $V(G)^k\to V(H)^k$: for a proposed image tuple, invert the
bijection at the first slot, then the bijection determined by that
prefix at the second slot, and continue.  Every matched pair of tuples
has equal colors at every round by the preceding induction.  Their
histograms therefore agree in either convention.
\end{proof}

The joint convention and its $(k+1)$-pair characterization are standard;
see, for example, \cite[Lemma~2]{CFICompression} and
\cite[arXiv v2, Lemma~7]{WLDimension}.  The proof above fixes explicitly the
infinite-game convention needed here.  A strategy can also be used in
reverse: announce the inverse bijection and translate Spoiler's selected
vertex through it.  Thus restricting Spoiler to select vertices in the
first graph does not change the winning relation.

\begin{lemma}[Pure tuples in a disjoint union]
\label{game:union}
Suppose $k\geq3$ and $G,H$ are connected graphs of diameter at most two.
Run either convention on $D=G\sqcup H$, including mixed tuples.  The
stable color histograms restricted to the pure tuples $V(G)^k$ and
$V(H)^k$ are equal if and only if the separate-input histograms from
Definition~\ref{game:conventions} are equal.
\end{lemma}

\begin{proof}
Stable colors in $D$ determine, for each two coordinates $u,v$, whether
\[
 u=v\quad\text{or}\quad E(u,v)\quad\text{or}\quad
 \exists z\,\bigl(E(u,z)\wedge E(z,v)\bigr).
\]
For the existential clause, replace a third coordinate and inspect its
two adjacencies in the atomic information retained by the replacement
colors.  This works for both update conventions.  By the diameter
assumption, the displayed condition is exactly membership in the same
component of $D$.

Take two related pure tuples, one in each component.  A replacement
leaves another coordinate as an anchor.  Consequently a stable-color
witness bijection on $D$, whether joint or for a specified slot, maps
replacement vertices in $G$ into $H$ and replacement vertices in $H$
into $G$.  Its restriction is a bijection $V(G)\to V(H)$.  Equal pure
histograms supply initial related tuples, and these restricted
bijections give exactly the invariant proof of
Lemma~\ref{game:wl}.  Hence Duplicator wins between $G$ and $H$.

Conversely, let Duplicator win between $G$ and $H$.  Play on $D$ against
itself by swapping the two components: use that strategy on vertices
selected in $G$ and its reverse on vertices selected in $H$.  After
any discards, each component has at most $K-1$ pairs when a new pair is
requested.  Both local strategies may therefore propose a bijection;
announce their union and commit the move only in the selected
component.  Cross-component entries have neither equality nor
adjacency on either side, so this is a winning strategy on $D$.
The successive-bijection argument in Lemma~\ref{game:wl} now matches
$V(G)^k$ bijectively with $V(H)^k$ while preserving stable colors in
$D$.  This proves equality of the pure histograms.
\end{proof}

\begin{remark}[Parameter used by the construction]
\label{game:parameter}
The remainder of the proof uses $K=r+1$ pairs and assumes $r\geq3$.
Under the joint convention $r=k$; under the separate convention
$r=k-1$.  Thus a query occupying $r$ pairs has exactly one spare pair
for transferring it.  The final graphs have a universal hub, so
Lemma~\ref{game:union} also applies to the convention that compares
pure tuples inside their disjoint union.
\end{remark}
